    En este proyecto se ha llevado a cabo el desarrollo de Colivi, una plataforma web orientada a resolver los problemas comunes en la busqeuda de alojamiento compartido y la gestión de la convivencia en pisos de estudiantes.

    Para solucionar esto, se ha implementado una arquitectura basada en un backend con Spring boot (java 21) y PostgreSQL, junto con un frontend interactivo desarrollado en React. Además, se ha integrado un microservicio en Node.js que implementa el protocolo MCP (Model Context Protocol) para incorporar un asistente conversacional capaz de consultar el estado de la vivienda, recomendar anuncios o resumir la actividad mediante modelos de lenguaje.

    El objetivo era consolidar buenas prácticas de ingeniería de software mediante principios SOLID y diseño guiado por el dominio (DDD), aprendiendo a resolver problemas reales de concurrencia, sincronización de datos y gamificación de tareas dando como resultado un sistema integral que cubre desde el primer contacto entre inquilino y propietario hasta el día a día dentro del hogar.

\vspace{1.5cm}
%\section*{Abstract}
    This project presents the design and implementation of Colivi, a web platform created to address common challenges in finding shared accommodations and managing household life in student flats. The main motivation behind the project stems from my personal experience during an Erasmus mobility program, where I faced significant hurdles searching for rooms from abroad, dealing with uncoordinated landlords, and managing daily domestic chores and living friction with flatmates.

    To solve these issues, the system relies on a backend built with Spring Boot (Java 21) and PostgreSQL, alongside a responsive frontend developed with React. Additionally, an external Node.js microservice implementing the Model Context Protocol (MCP) was integrated to power an intelligent assistant capable of querying household chores, recommending accommodation options, and summarizing inbox activity using language models.

    The primary goal was to strengthen software engineering best practices through SOLID principles and domain-driven design, learning to tackle real-world concurrency challenges, data synchronization, and gamification mechanisms. The final outcome is a comprehensive platform that supports users throughout the entire flat-sharing lifecycle, from initial room discovery to day-to-day household organization.
